\chapter{UI/UX Enhancements}
\label{ch:ui}

\emph{Pull requests: \#189, \#192, \#246.}

\section{Overview}
This chapter gathers three user-experience improvements that sharpened the OsdagBridge
interface: a custom pointing-hand cursor applied consistently across the application
(PR~\#189), a \emph{Legends} toggle in the toolbar together with a cleaner axis gizmo
(PR~\#192), and conditional visibility for the optional Railing and Median components so the
CAD panel only offers toggles for parts the design actually contains (PR~\#246).

\section{Custom Cursor}
Qt's built-in pointing-hand cursor looks different across platforms, and on Linux in particular
it did not match the rest of the UI. PR~\#189 introduced a programmatically drawn cursor in a
new module, \texttt{desktop/ui/utils/custom\_cursors.py}. The hand is defined as a small pixel
grid and rendered at 8$\times$ supersampled resolution, then scaled down for smooth edges:

\begin{lstlisting}[caption={Programmatically generated pointing-hand cursor (\texttt{custom\_cursors.py}).},label={lst:cursor}]
def _create_pointing_hand_pixmap(size: int = 32) -> QPixmap:
    """Create a classic upright pointing-hand cursor with smooth edges."""
    cursor_data = [                 # 0 = transparent, 1 = outline, 2 = fill
        "00000000000111100000000000000000",
        "00000000001222100000000000000000",
        # ... 32 rows total ...
    ]
    supersample = 8                 # render hi-res, then scale down smoothly
    ...

def pointing_hand_cursor() -> QCursor:
    """Custom hand cursor on Linux, standard cursor elsewhere."""
    return get_cursor(Qt.CursorShape.PointingHandCursor)
\end{lstlisting}

Every place that previously used \texttt{Qt.PointingHandCursor} now calls
\texttt{pointing\_hand\_cursor()} instead---the checkboxes in \texttt{cad\_3d.py}, the additional
-inputs button in \texttt{input\_dock.py}, dialog buttons in \texttt{material\_properties.py},
and the base custom button in \texttt{custom\_buttons.py}---so the whole application shares one
consistent cursor.

\section{Legend Toggle and a Cleaner Axis}
PR~\#192 added a \emph{Legends} button to the toolbar. Like the other toggles, the button is
created in \texttt{custom\_widgets.py} and wired by the controller; toggling it shows or hides
the bottom-left information table (\texttt{\_info\_table}) that lists the bridge parameters:

\begin{lstlisting}[caption={The legend toggle button and its handler (\texttt{custom\_widgets.py}, \texttt{toolbar\_controller.py}).},label={lst:legend}]
# custom_widgets.py -- create the button
create_button(":/vectors/tool_bar/legend_icon.svg", "Legends")

# toolbar_controller.py -- find and toggle it
_TIP_LEGEND = "Legends"
self._btn_legend = self._find_button(self._TIP_LEGEND)

def _cad_toggle_legend(self):
    table = self._cad_widget._info_table
    if self._btn_legend.isChecked():
        table.show()
        self._cad_widget._position_info_table()
    else:
        table.hide()
\end{lstlisting}

The same PR also removed the grey rounded box that had been drawn behind the coordinate-axis
triad (Chapter~\ref{ch:axis}). The triad's \texttt{paintEvent} now draws only the axis arrows on
a translucent background, dropping the \texttt{drawRoundedRect} call and the grey brush, so the
gizmo sits cleanly over the model.

\section{Optional-Component Visibility}
The CAD component panel offers a checkbox per bridge part. But not every design has a footpath
railing or a median---showing those toggles regardless was misleading. PR~\#246 made Railing and
Median \emph{optional}: their toggles appear only when the design actually includes them.

A set marks which components are optional, and \texttt{set\_available\_components} hides the ones
the current design does not have, remembering them in \texttt{self.\_unavailable}:

\begin{lstlisting}[caption={Marking components optional and hiding unavailable ones (\texttt{cad\_3d.py}).},label={lst:optional}]
OPTIONAL_KEYS = {"Railing", "Median"}

# availability from the design parameters
railing_available = getattr(design_params, "footpath_config", "NONE") != "NONE"
median_available  = getattr(design_params, "enable_median", False)

def set_available_components(self, available):
    """Show toggles for available components; uncheck + hide the rest."""
    for key in OPTIONAL_KEYS:
        if key not in available:
            self._checkboxes[key].setChecked(False)
            self._checkboxes[key].hide()
            self._unavailable.add(key)
\end{lstlisting}

Crucially, hidden components are also excluded when the panel \emph{collects} the selected
components for rendering, so an absent part can never be switched on---even through the master
``Bridge'' toggle:

\begin{lstlisting}[caption={Excluding unavailable components from the render selection (\texttt{cad\_3d.py}).},label={lst:collect}]
def _collect(self):
    selected = []
    for key, cb in self._checkboxes.items():
        if key in self._unavailable:
            continue                  # never render a component the design lacks
        if cb.isChecked():
            selected.append(key)
    return selected
\end{lstlisting}

Availability is recomputed on every render (in \texttt{render\_3d\_cad}, before the model is
loaded), so the toggles correctly reappear when the user switches to a design that does contain
a footpath or median.

\section{Result}
The application now shares one consistent hand cursor, offers a Legends toggle over a cleaner
axis gizmo, and shows Railing/Median toggles only when the design includes those components---so
the CAD panel always matches the actual geometry.

%% IMAGE NEEDED: custom_cursor.png -- the custom pointing-hand cursor over the UI
\osdagfig{custom_cursor.png}{The custom pointing-hand cursor used consistently across
OsdagBridge.}{fig:cursor}

%% IMAGE NEEDED: optional_components.png -- CAD panel where Railing/Median toggles are hidden vs shown
\osdagfig{optional_components.png}{The Railing and Median toggles appear only when the design
includes those components.}{fig:optional}

\section{Summary}
These three enhancements---a consistent custom cursor, a legend toggle with a cleaner axis, and
design-aware optional-component visibility---made the OsdagBridge interface more consistent and
less confusing, ensuring the controls the user sees always reflect the current design.
